\documentclass[11pt,xcolor={dvipsnames},aspectratio=159,hyperref={pdftex,pdfpagemode=UseNone,hidelinks,pdfdisplaydoctitle=true},usepdftitle=false]{beamer}
\usepackage{presentation}[aspectratio=169]
\usepackage{math}
\usepackage{mathtools}
\usepackage{amsmath}
\usepackage{mleftright}

% Enter title of presentation PDF:
\hypersetup{pdftitle={Price level}}
% Enter link to PDF file with figures:


\begin{document}
% Enter presentation title:
\title{Price level determination in a monetary economy}
\subtitle{Fiscal and Monetary Policy 2024}
% Enter presentation information:

% Enter presentation authors:
\author{Piotr Żoch}%
% Enter presentation location and date (optional; comment line if not needed):
\frame{\titlepage}

% Fill out content of presentation:
\begin{frame}{Plan}
    \begin{itemize}
    \item We study how the price level (and inflation) is determined in a monetary economy.
    \item This will later serve us as a framework to study interactions of fiscal and monetary policy.
    \item We will start with a very simple environment and then gradually add more features.
    \end{itemize}
\end{frame}


\begin{frame}
    \heading{Endowment economy}
\end{frame}

\begin{frame}{Endowment economy}
    \begin{itemize}
    \item Every period there is an endowment of nonstorable goods $y_t$.
    \item For simplicity assume $y_t = y$ for all $t$.
    \item Let $P_t$ be the price level in period $t$: one unit of goods costs $P_t$ units of account.
    \item There is a representative household that trades contingent claims (Arrow securities).
    \item The household maximizes $$\E \sum_{t=0}^\infty \beta^t u\of{c_t}.$$
    \item No government purchases of goods. 
    \end{itemize}
\end{frame}

\begin{frame}{Asset prices}
\begin{itemize}
        \item Let $q_{t,t+1}$ be the real price (in goods) of a claim that pays one unit of goods in period $t+1$ in a particular state of the world.
        \item The nominal contingent claim price is \begin{align*}
           Q_{t,t+1} = q_{t,t+1} \frac{P_t}{P_{t+1}}.
        \end{align*}
        \item The \al{nominal interest rate} $i_t$ satisfies \begin{align*}
            \frac{1}{1+i_t} = \E_t Q_{t,t+1}.
        \end{align*}
    \end{itemize}
\end{frame}

\begin{frame}{Asset prices}
\begin{itemize}
        \item In this economy $c_t = y$.
        \item The (real) SDF, $\beta u'\of{c_t+1}\slash u'\of{c_t}$, is constant and equal to $\beta$.
        \item The \al{real interest rate}  $r$ is also constant and satisfies \begin{align*}
            \frac{1}{1+r} = \beta.
        \end{align*}
        \item The \al{nominal discount factor} is $Q_{t,t+1} = \beta \frac{P_t}{P_{t+1}}$. 
    \end{itemize}
\end{frame}

\begin{frame}{Asset prices}
\begin{itemize}
        \item We define gross inflation rate as 
         \begin{align*}
            \Pi_{t+1} \coloneq \frac{P_{t+1}}{P_t}.
        \end{align*}
        \item We obtain the \alb{Fisher equation}
        \begin{align*}
            \frac{1}{1+i_t} = \frac{1}{1+r} E_t \frac{P_{t}}{P_{t+1}} = \frac{1}{1+r} E_t\frac{1}{\Pi_{t+1}} = \beta  E_t\frac{1}{\Pi_{t+1}} .
         \end{align*}
        \item The usual formulation of the Fisher equation is obtained by linearization: 
        \begin{align*}
                i_t = r + \E_t \pi_{t+1}.
        \end{align*}
    \end{itemize}
\end{frame}


\begin{frame}{Fisher equation}
    \begin{itemize}
        
            \item In general, the Fisher equation is \begin{align*}
                i_t = r_t + \E_t \pi_{t+1}.
            \end{align*}
            \item This is a no-arbitrage condition: you should be indifferent between two types of investment (real and nominal).
        \end{itemize}
    \end{frame}

\begin{frame}{Determinacy}
    \begin{itemize}
            \item In this endowment economy \begin{align*}
                \frac{1}{1+i_t} = \beta  E_t\frac{1}{\Pi_{t+1}}
            \end{align*}
            is the \alb{only} equilibrium condition -- we already used goods market clearing.
            \item What does it imply for the sequence of price levels $\bc{P_t}_{t=0}^\infty$?
            \item In particular, we say that the level of inflation is \al{unique} or \al{determinate} in equilibrium if: \begin{enumerate}
                \item There is a \al{unique} scalar $P_0$ in equilibrium.
                \item If $\Pi^\prime$ and $\Pi^{\prime\prime}$ are two sequences that satisfy equilibrium conditions, then $\Pi^{\prime} = \Pi^{\prime\prime}$.
            \end{enumerate}
        \end{itemize}
    \end{frame}

\begin{frame}{Determinacy}
    \begin{itemize}
            \item The question of determinacy may seem esoteric, but it is important.
            \item It boils down to the question of whether \al{some policy can control prices}.
            \item In this simple economy, the only available policy is the choice of the nominal interest rate.
            \item Suppose we have some target for the price level (or inflation rate) -- can we achieve it by setting interest rates?
        \end{itemize}
    \end{frame}

\begin{frame}{Determinacy}
    \begin{itemize}
            \item Consider the simplest possible policy: \al{interest rate peg}.
            \item The central bank sets $i_t = i$ for all $t$ -- it does not say it will \emph{respond} to something, it is \al{committed} to $i$ at all times in all states of the world.
            \item We have \begin{align*}
                \frac{1}{1+i} = \beta  E_t\frac{1}{\Pi_{t+1}}.
            \end{align*}
            \item This determines $E_t\frac{1}{\Pi_{t+1}}$  -- approximately the \alb{expected} inflation rate.
            \item But it does \alr{not} determine $\Pi_{t+1}$ -- the \alb{actual} inflation rate.  
            \item Agents expect $1\slash\Pi_t+1$ to be equal to $1/(1+i)$, but the actual inflation rate is \begin{align*}
            \frac{1}{\Pi_{t+1}} = E_t\frac{1}{\Pi_{t+1}} + \epsilon_{t+1},\quad E_t\epsilon_{t+1} = 0. 
        \end{align*}
   
        \end{itemize}
    \end{frame}

    \begin{frame}{Determinacy}
        \begin{itemize}
                \item What if we have \al{perfect foresight}? 
                \item Then it is not a problem, because $\epsilon_{t+1} = 0$. 
                \item For all $t\geq0$ $\Pi_{t+1}$ is determined by the peg, but $P_0$ is not (because $\Pi_0$ is not).
                \item This is because we do not have a Fisher equation for $t=0$.
                \item Conclusion: interest rate peg \alr{does not} ensure price level determinacy. 
            \end{itemize}
        \end{frame}

    \begin{frame}{Determinacy}
        \begin{itemize}
                \item To simplify the analysis, we will now focus on \al{linearized} Fisher equation \begin{align*}
                    i_t = r + \E_t \pi_{t+1}.
                \end{align*}
                \item Consider an \al{interest rate rule} $$i_t =  r + \phi \pi_t$$.
                \item We call these rules \al{Taylor rules}. In general there can be other variables in the rule, and the intercept can be different than $r$.
                \item Use the expression for $i_t$ in the Fisher equation to obtain \begin{align*}
                 \phi \pi_t = \E_t \pi_{t+1}. 
                \end{align*}
            \end{itemize}
        \end{frame}

    \begin{frame}{Determinacy}
        \begin{itemize}
                \item We have the difference equation \begin{align*}
                    \pi_t = \frac{1}{\phi}\E_t \pi_{t+1}. 
                   \end{align*}
                \item We can solve it \al{forward} as 
                   \begin{align*}
                    \pi_t = \lim_{T\to\infty} \phi^{-T} \E_t \pi_{t+T}.
                   \end{align*} 
                \item If $\phi>1$ \al{and} $E_t \pi_{t+T}$ \alr{does not diverge}, then $\pi_t=0$ for all $t$.
            \end{itemize}
        \end{frame}

    \begin{frame}{Determinacy}
        \begin{itemize}
                \item If $\phi>1$ \al{and} $E_t \pi_{t+T}$ \alr{does not diverge}, inflation is determinate in all periods, even at $t=0$.
                \item $\pi_0$ is determined, so $P_0$ is determined too. 
                \item If the response of the interest rate to inflation is \al{strong enough}, $\phi>0$ then $\pi_t = 0$ for all $t$ is the only bounded solution.
                \item We call $\phi>1$ the \al{Taylor principle}.
                \item If $\phi<1$, then $\pi_t$ is \al{not} determined: there are multiple bounded solutions of the difference equation.
                \item Where does $E_t \pi_{t+T}$ not diverging come from? 
            \end{itemize}
        \end{frame}

    \begin{frame}{Determinacy}
        \begin{itemize}
                \item The terminal condition on $E_t \pi_{t+T}$ is \alr{not} an optimality condition. 
                \item Even if $\phi>1$, but if the household for whatever reason expects inflation to diverge (i.e., hyperinflation), it will.
                \item In this model $\phi>1$ has nothing to do with the logic like "increase real rates to reduce aggregate demand"
                \item In a way, the central bank makes a threat: if you do not behave, I will make inflation explode. Is it reasonable? Is it how central banks do it?
                \item See Cochrane (2011) for a discussion of Taylor rules.
                \item My advice: you might not like the model (it is obviously simple) -- but do not apply some outside logic to understand these results!
            \end{itemize}
        \end{frame}

        \begin{frame}{Determinacy}
            \begin{itemize}
                    \item The usual justification is that when inflation gets out of control, the central bank will switch to something else.
                    \item But then, it is not a Taylor rule with $\phi>1$ that ensures determinacy!
                    \item What is this other thing that would ensure determinacy?
                    \item For example, the ECB has a monetary pillar, understood as a commitment to switch to a monetary approach to pin down inflation if inflation starts exploding.
                    \item We will soon turn to these monetary approaches.
                \end{itemize}
            \end{frame}


        \begin{frame}{Determinacy}
        \begin{itemize}
                \item We used $r$ as the intercept. Suppose $r_t$ is not constant and the central bank follow a more standard Taylor rule: 
                \begin{align*}
                    i_t = \bar{r}_t + \pi^* + \phi \bp{\pi_t-\pi^*},
                \end{align*}
                where $\bar{r}_t$ is the intercept, $\pi^*$ is the target inflation rate and $\phi$ is the response to inflation.
                \item If we use this rule in the Fisher equation, we obtain \begin{align*}
                   \pi_{t}-\pi^{*} & =\frac{1}{\phi}\left[r_{t}-\bar{r}_t+\left(E_{t}\left[\pi_{t+1}\right]-\pi^{*}\right)\right]
                \end{align*}
                \item What is needed to ensure that $\pi_t = \pi^*$ for all $t$?
            \end{itemize}
        \end{frame}



        \begin{frame}{Determinacy}
            \begin{itemize}
                    \item The problem is now harder, because $r_t$ possibly fluctuates.
                    \item In a richer model $r_t$ would depend on the monetary policy itself: monetary policy would affect the real economy. 
                    \item To know how, you would need to know how inflation is linked to output (e.g., through the Phillips curve).
                    \item Holden (2024) shows that if you replace $\bar{r}_t$ with $r_t$ in the monetary policy rule, you will restore determinacy -- \alb{Robust Real Rate Rules} (subject to the qualification that $E_t \pi_{t+T}$ does not diverge).
                    \item If there is a unit response to real interest rates, i.e., $\bar{r} = r_t$, then you do not need to know other details of the economy to set inflation at the target. You \al{only} need a Fisher equation in the economy and you need to observe $r_t$ somehow. 
                    \item This is an \al{extremely} powerful result. 
                \end{itemize}
            \end{frame}


        \begin{frame}{Determinacy}
            \begin{itemize}
                    \item Assume there is some money demand in this economy: \begin{align*}
                    m_t = p_t +  y - \eta i_t,\end{align*} where $\eta>0$ is a semielasticity of money demand.
                    \item Use the Fisher equation to obtain 
                    \begin{align*}
                        m_t = p_t +  y - \eta \bs{r + \E_t \pi_{t+1}}.
                    \end{align*}
                    \item Let the central bank follow some money growth rule, e.g. $m_t =  m_{t-1} + \bar{m} $.
                    \item Is it enough to ensure price level determinacy?
                \end{itemize}
            \end{frame}

            
        \begin{frame}{Determinacy}
            \begin{itemize}
                    \item We have a formula \begin{align*}
                        \bp{1+\eta} \bp{ p_t  - \bar{m} \cdot t  } =\eta \bp{\E_t p_{t+1} - \bar{m} \cdot t} - y + p_t + \eta r.
                    \end{align*}
                    we can solve it forward and there will be a term like $\lim_{T\rightarrow\infty} \E_t p_{t+T} - \bar{m} \cdot \bp{t+T-1} $.
                    \item This term will be equal to 0 by the household's optimality condition (transversality condition).
                    \item \al{Intuition}: households do not want to hold too much or too little money. 
                    \item We get \begin{align*}
                        p_t  = \bar{m}\cdot t + \eta \bar{m} + \frac{1}{1+\eta} \sum_{j=0}^\infty \bp{\frac{\eta}{1+\eta}}^j \bp{\eta r - y}.
                    \end{align*}
                \end{itemize}
            \end{frame}

            \begin{frame}{Determinacy}
                \begin{itemize}
                        \item Money growth rule ensures price level determinacy (in this case).
                        \item This is true even if there are shocks $u_t$ to money demand.
                        \item With different microfoundations for money demand, conclusions might be different.
                        \item We can solve for the implied nominal interest rate: \begin{align*}
                            i_t = \frac{1}{\eta}\bp{p_t - m_t + y}
                        \end{align*}
                        \item Note it responds to \al{price level}. This is an equilibrium outcome. 
                        \item If we instead \al{started} with a rule that responds to price levels it would be a \alb{Wicksellian rule}.
                \end{itemize}
                \end{frame}


\begin{frame}{Budget constraint?}
    \begin{itemize}
            \item So far we have seen that fiscal policy can affect inflation / price levels by affecting supply / demand of money. 
            \item Is there any other link? In particular, what if the economy is \al{cashless}?
            \item One possibility is that fiscal policy has effects on output and then on inflation through the Phillips curve.
            \item We will focus on a different link that could be present even in an endowment economy.
            \item That link will depend on how we interpret the government budget "constraint".
        \end{itemize}
    \end{frame}

\begin{frame}{Budget constraint?}
    \begin{itemize}
            \item Example: simple static economy.
            \item In the morning, bondholders wake up and own $B_0$ units of zero coupon bonds that promise to pay one unit of account ("dollar") each. 
            \item The government pays bondholders by printing ``money".
            \item In the evening, the government collects taxes $P_1 T_1$. Taxes are paid in "money". 
            \item The world ends aferwards.
        \end{itemize}
    \end{frame}

\begin{frame}{Budget constraint?}
    \begin{itemize}
            \item In equilibrium \begin{align*}
                B_0 = P_1 T_1.
            \end{align*}
            \item  ``Money" printed up in the morning must be all ``soaked" up by taxes at the end of the day.
            \item $B_0$ is predetermined. We assume $T_1$ is exogenous. 
            \item The price level $P_1$\al{adjusts} to ensure this equality.
            \item This is an \alb{equilibrium condition}.
            \item It looks a bit like a budget constraint, but it is \alr{not}.
        \end{itemize}
    \end{frame}

\begin{frame}{Budget constraint?}
    \begin{itemize}
            \item In the example above we did not require the government to "collect enough taxes to pay its debt".
            \item This is the \alr{Fiscal Theory of the Price Level}.
            \item Consider other possible stories about how $P_1$ adjusts to ensure $B_0 = P_1 T_1$.
            \item If $P_1$ is too low, more money was printed that will be collected in taxes. People will try to get rid of money by buying goods. This will push up the price level. 
            \item ``Too much money chasing too few goods".
            \item Too much government debt relative to fiscal surpluses is like net wealth -- people will try to spend more. 
        \end{itemize}
    \end{frame}
  
    
\begin{frame}{FTPL intro}
    \begin{itemize}
            \item Contrast this with the \alb{quantity theory of money}.
            \item In the quantity theory of money, inflation results from more money than needed to settle transations or to satisfy other money demand for other reasons. All money supply matters. 
            \item Here only \alg{outside money} -- government liabilities -- matter. It is also irrelevant whether it is debt or currency. 
            \item In the example it is not essential that the government "prints" money: people could pay taxes directly with bonds. 
        \end{itemize}
    \end{frame}

\begin{frame}{FTPL intro}
    \begin{itemize}
            \item The government issues the currency and nominal debt that define the price level.
            \item One possible interpretation nominal government debt is like corporate \emph{equity}: its price adjusts to make a valuation equation hold.
            \item This does not hold for real or foreign currency government debt. 
            \item This does not hold for personal or corporate debt.
        \end{itemize}
    \end{frame}
\begin{frame}{FTPL intro}
    \begin{itemize}
            \item Note that the government could \al{choose} to set taxes $T_1$ such that $T_1 = \frac{B_0}{P_1}$ for all possible $P_1$, this means $T_1 = f\of{P_1}$.
            \item In this case the price level is \alb{indeterminate}.
            \item We call such a fiscal policy \al{passive} or \al{Ricardian}.
            \item The point is that the government does not have to follow such a policy. Fiscal policy might be \al{active} or \al{non-Ricardian}. 
        \end{itemize}
    \end{frame}

    \begin{frame}{An intertemporal model}
        \begin{itemize}
                \item We now consider the most basic fully-specified version of the FTPL.
                \item A representative households maximizes \begin{align*}
                    \sum_{t=0}^\infty \beta^t u\of{c_t}
                \end{align*}
                in a complete asset market economy, with restriction $B_t \geq 0$ and $M_t \geq 0$.
                \item The household has a constant endowment every period, $y$
                \item The budget constraint of the household is \begin{align*}
                    M_{t-1} + B_{t-1} + P_t y = P_t c_t + P_t s_t +  M_t + Q_t B_t \end{align*}
                    where $Q_t \coloneq \frac{1}{1+i_t}$ is the price of bonds in terms of money and $s_t$ is real taxes net of transfers. 
            \end{itemize}
        \end{frame}

\begin{frame}{An intertemporal model}
    \begin{itemize}
            \item Notice that if $Q_t < 1$, then the household does not hold any money, $M_t = 0$.
            \item If $Q_t = 1$, then money and bonds are perfect substitutes and we can use $B_t$ to denote their sum.
            \item $Q_t > 1$ not possible. 
            \item Drop $M_t$ from the budget constraint and use market clearing $c_t = y$ to obtain \begin{align*}
                B_{t-1}  = P_t s_t +  Q_t B_t \end{align*}
        \end{itemize}
    \end{frame}

\begin{frame}{An intertemporal model}
    \begin{itemize}
            \item Household optimality together with the market clearing implies \begin{align*}
                Q_t = \beta \E_t \frac{P_t}{P_{t+1}}.
            \end{align*}

            \item There is also the transversality condition 
            \begin{align*}
                \lim_{T\to\infty} \beta^T \E_t \frac{B_{T-1}}{P_T}  = 0.
            \end{align*}
            \item We obtain \begin{align*}
                \frac{B_{t-1}}{P_t}  = \E_t \sum_{j=0}^\infty \beta^j s_{t+j}.
                \end{align*}
        \end{itemize}
    \end{frame}

\begin{frame}{An intertemporal model}
    \begin{itemize}
            \item The government debt valuation equation \begin{align*}
                \frac{B_{t-1}}{P_t}  = \E_t \sum_{j=0}^\infty \beta^j s_{t+j}.
                \end{align*}
                determines the price level $P_t$.
            \item It does so in every period. 
            \item This looks like a standard asset pricing formula! 
            \item The government debt looks like a stock. Price per share $1/P_t$ times number of shares $B_{t-1}$ equals the present value of dividends $\bc{s_{t+j}}$.
            \item There are problems with this analogy, it is potentially misleading -- see Hagedorn (2024). 
        \end{itemize}
    \end{frame}

\begin{frame}{An intertemporal model}
    \begin{itemize}
            \item What if the government changes nominal debt without changing surpluses?
            \item What if the government changes surpluses without changing nominal debt?
            \item It is useful to rewrite the equation as \begin{align*}
                \frac{B_{t}}{P_{t+1}}  = \E_{t+1} \sum_{j=0}^\infty \beta^j s_{t+j+1}.
                \end{align*}
            \item Define $\Delta\E_{t+1} \coloneq \E_{t+1} - \E_{t}$ to write \begin{align*}
                \frac{B_{t}}{P_{t}} \Delta\E_{t+1}\frac{P_{t+1}}{P_t}  = \Delta\E_{t+1} \sum_{j=0}^\infty \beta^j s_{t+j+1}.
                \end{align*}
        \end{itemize}
    \end{frame}

\begin{frame}{An intertemporal model}
    \begin{itemize}
            \item At time $t+1$ both $B_t$ and $P_t$ are predetermined. 
            \item \al{Unexpected} inflation is determined by changing expectations of the present value of surpluses.
            \item Bad fiscal news affect inflation for one period only. 
            \item This is because we assumed one-period debt only.
        \end{itemize}
    \end{frame}

\begin{frame}{An intertemporal model}
    \begin{itemize}
            \item We can also rewrite the equation as \begin{align*}
                \frac{B_{t}}{P_{t}} \E_{t}\frac{P_{t+1}}{P_t}  = \E_{t} \sum_{j=0}^\infty \beta^j s_{t+j+1}.
                \end{align*}
            \item It is the same as \begin{align*}
                \frac{B_{t}}{P_{t}} \frac{1}{1+i_t}  = \E_{t} \sum_{j=0}^\infty \beta^j s_{t+j+1}.
                \end{align*}
            \item The left hand side is the real revenue the government gets from selling debt $B_t$ at the end of period $t$.
            \item The price level $P_t$ is already determined at the beginning of period $t$.
            \item The revenue from bonds sales is fully determined by the PDV of surpluses. 
        \end{itemize}
    \end{frame}    

\begin{frame}{An intertemporal model}
    \begin{itemize}
            \item The government can target nominal interest rates by changing $B_t$ without affecting surpluses. 
            \item By standard arguments the expected inflation rate depends on the nominal interest rate. 
            \item Monetary policy: changing $i_t$ / changing $B_t$ without adjusting surpluses. 
            \item Fiscal policy: changing surpluses without adjusting $i_t$ / $B_t$.
        \end{itemize}
    \end{frame}    

\end{document}